\subsection{Composizione, inversa e numerabilità}
\label{sec:esercizi-fondamenti-composizione}

\begin{exercise}[Inversa di una funzione affine]
\label{ex:fondamenti-inversa-funzione-affine}
\difficoltaesercizio{0}
Dimostrare che $f\colon\R\longrightarrow\R$, definita da
$f(x)=3x-5$, è biiettiva e determinarne l'applicazione inversa,
con dominio e codominio.
\rimandoteoria{def:lez06-iniettiva,def:lez06-suriettiva,def:lez09-inversa}
\end{exercise}
\begin{proof}[Soluzione dell'esercizio]
\emph{Iniettività.} Siano $x_1,x_2\in\R$. Allora
\[
  f(x_1)=f(x_2)
  \quad\Longrightarrow\quad
  3x_1-5=3x_2-5
  \quad\Longrightarrow\quad x_1=x_2.
\]

\emph{Suriettività.} Per ogni $y\in\R$, l'equazione
\[
  y=3x-5
  \quad\Longleftrightarrow\quad
  x=\frac{y+5}{3}
\]
ammette una soluzione in $\R$. Pertanto $f$ è biiettiva e
\[
  f^{-1}\colon\R\longrightarrow\R,
  \qquad f^{-1}(x)=\frac{x+5}{3}.
\]
\end{proof}

\begin{exercise}[Inversa e restrizione del dominio]
\label{ex:fondamenti-inversa-restrizione-dominio}
\difficoltaesercizio{0}
Per
\[
  f\colon[1,+\infty)\longrightarrow[0,+\infty),
  \qquad f(x)=(x-1)^2,
\]
verificare l'iniettività e la suriettività e scrivere
l'applicazione inversa.
\rimandoteoria{def:lez06-iniettiva,def:lez06-suriettiva,def:lez09-inversa}
\end{exercise}
\begin{proof}[Soluzione dell'esercizio]
\emph{Iniettività.} Per $x_1,x_2\ge1$,
\[
  (x_1-1)^2=(x_2-1)^2
  \quad\Longrightarrow\quad x_1=x_2,
\]
poiché $x_1-1,x_2-1\ge0$.

\emph{Suriettività.} Per ogni $y\ge0$, l'equazione
\[
  y=(x-1)^2
  \quad\Longleftrightarrow\quad x=1+\sqrt{y}
\]
ammette una soluzione in $[1,+\infty)$. Pertanto $f$ è
biiettiva e
\[
  f^{-1}\colon[0,+\infty)\longrightarrow[1,+\infty),
  \qquad f^{-1}(x)=1+\sqrt{x}.
\]
\end{proof}

\begin{exercise}[Ordine delle composizioni]
\label{ex:fondamenti-ordine-composizioni}
\difficoltaesercizio{0}
Siano $\varphi,\psi\colon\R\longrightarrow\R$ definite da
$\varphi(x)=x^2+1$ e $\psi(x)=2x-3$.
Calcolare $\varphi\circ\psi$ e $\psi\circ\varphi$ e trovare
un valore di $x$ in cui differiscono.
\rimandoteoria{def:lez09-composizione}
\end{exercise}
\begin{proof}[Soluzione dell'esercizio]
Per la definizione di composizione,
\[
\begin{aligned}
  (\varphi\circ\psi)(x)
    &=(2x-3)^2+1=4x^2-12x+10,\\
  (\psi\circ\varphi)(x)
    &=2(x^2+1)-3=2x^2-1.
\end{aligned}
\]
Per $x=2$ si ha
\[
  (\varphi\circ\psi)(2)=2,
  \qquad (\psi\circ\varphi)(2)=7.
\]
Pertanto le due composizioni non coincidono.
\end{proof}

\begin{exercise}[Una nuova classe di naturali numerabile]
\label{ex:fondamenti-classe-naturali-numerabile}
\difficoltaesercizio{1}
Sia
\[
  A=\{3n+2:n\in\N\}.
\]
Dimostrare che $A$ è numerabile costruendo una biezione esplicita
$\varphi\colon\N\longrightarrow A$ e verificando separatamente
iniettività e suriettività.
\rimandoteoria{def:lez07-numerabilita}
\end{exercise}
\begin{proof}[Soluzione dell'esercizio]
Si consideri l'applicazione
\[
  \varphi\colon\N\longrightarrow A,
  \qquad \varphi(n)=3n+2.
\]
\emph{Iniettività.} Per $n_1,n_2\in\N$,
\[
  \varphi(n_1)=\varphi(n_2)
  \quad\Longrightarrow\quad
  3n_1+2=3n_2+2
  \quad\Longrightarrow\quad n_1=n_2.
\]

\emph{Suriettività.} Per ogni $a\in A$ esiste $n\in\N$
tale che $a=3n+2=\varphi(n)$.
Pertanto $\varphi$ è biiettiva e $A$ è numerabile.
\end{proof}

\clearpage

\begin{exercise}[Inversa di un rapporto razionale]
\label{ex:fondamenti-inversa-rapporto-razionale}
\difficoltaesercizio{1}
Per
\[
  f\colon\R\setminus\{-1\}\longrightarrow\R\setminus\{2\},
  \qquad f(x)=\frac{2x+3}{x+1},
\]
dimostrare che $f$ è biiettiva, calcolare $f^{-1}$ e
verificare entrambe le identità di composizione con l'identità.
\rimandoteoria{def:lez09-inversa,prop:lez09-inversa-composizione}
\end{exercise}
\begin{proof}[Soluzione dell'esercizio]
Per ogni $x\ne-1$ si ha
\[
  f(x)=\frac{2x+3}{x+1}=2+\frac{1}{x+1}.
\]
\emph{Iniettività.} Per $x_1,x_2\ne-1$,
\[
  f(x_1)=f(x_2)
  \quad\Longrightarrow\quad
  \frac1{x_1+1}=\frac1{x_2+1}
  \quad\Longrightarrow\quad x_1=x_2.
\]

\emph{Suriettività.} Per ogni $y\ne2$,
\[
  y=2+\frac1{x+1}
  \quad\Longleftrightarrow\quad
  x=\frac1{y-2}-1.
\]
Il valore ottenuto appartiene al dominio, poiché è diverso da $-1$.
Pertanto $f$ è biiettiva e
\[
  f^{-1}\colon\R\setminus\{2\}\longrightarrow
    \R\setminus\{-1\},
  \qquad f^{-1}(x)=\frac1{x-2}-1.
\]

Infine, per ogni $x\ne-1$ e ogni $y\ne2$, si ha
\[
\begin{aligned}
  (f^{-1}\circ f)(x)
    &=\frac1{f(x)-2}-1
     =\frac1{\frac1{x+1}}-1=x,\\
  (f\circ f^{-1})(y)
    &=2+\frac1{f^{-1}(y)+1}
     =2+\frac1{\frac1{y-2}}=y.
\end{aligned}
\]
\end{proof}
